\subsection{I-45. Quadratic residue constraints strengthen the reset-Hessian argument}
\label{I-45}

\paragraph{Statement.}
Fix distinct nodes over a field $E$ and retain the balanced
rational two-pole map with common center. Every nonzero
homogeneous multilinear common-translation-invariant
polynomial of ROABP width at most three in some order is
hit if
\[
 0\le n-2\deg f\le7,
 \quad\hbox{or if}\quad
 0\le n-2\deg f\le9\ \hbox{and }\operatorname{char}E\ne2.
 \tag{I.187}
\]
Consequently all nonzero multilinear width-at-most-three
polynomials on at most fifteen variables are hit in every
characteristic, and those on at most seventeen variables
are hit outside characteristic two. No order information
is supplied to the evaluator.

The finite hitting set for total degree at most $D$ is still
\[
 (n+1)^2\binom{D+3}{3},\qquad
 |E|\ge\max\{2n+1,D+1\},
 \tag{I.188}
\]
with the working-extension convention of I-40. These
assertions concern the balanced rational map, not an
unweighted ordinary centered SV map.

The two new ingredients are proved separately below.
Actual quadratic Taylor terms satisfy a one-pole residue
constraint that the abstract sharp Hessians of I-42 were
not required to satisfy. The constrained connector proof
of I-44 also extends from five to seven intervening cores.
Neither ingredient assumes preservation against arbitrary
parameter-dependent matrix boundaries.

\paragraph{1. The quadratic Taylor term belongs to the one-pole ideal.}
Suppose a multilinear polynomial $f$ annihilates balanced
two-pole evaluation. Over
$K=\overline E(t,z,y_0)$, put
\[
 \alpha_i=t+\frac{z}{a_i-y_0},\qquad
 h_i=\frac1{a_i-y_0},\qquad
 H_{ij}=(\partial_i\partial_j f)(\alpha)\quad(i\ne j),
 \quad H_{ii}=0.
 \tag{I.189}
\]
Pole extraction gives $(\partial_i f)(\alpha)=0$.
Differentiating these identities in $t,z,y_0$ yields the
three previously used kernels
\[
 H\mathbf1=H\mathbf h=H\mathbf h^{\circ2}=0.
 \tag{I.190}
\]
Here $z$ is a nonzero element of the rational function
field, so it can be divided out in the last identity.

There is a further constraint in every characteristic.
The degree-two Taylor polynomial, defined without a
factorial, is
\[
 Q_H(X)=\sum_{i<j}H_{ij}X_iX_j.
 \tag{I.191}
\]
For fresh independent $r,s,u,\varepsilon$, the assumed
two-pole identity gives
\[
 f\!\left(\alpha_i+
       \varepsilon\left(r+\frac{s}{a_i-u}\right)\right)_{i=1}^n=0.
\]
Take the coefficient of $\varepsilon^2$. Multilinearity
makes it exactly
$Q_H(r\mathbf1+s(1/(a_i-u))_i)=0$. Thus $Q_H$ belongs
to the balanced one-pole ideal over $K$. Its simple pole
at $u=a_i$ gives
\[
 \sum_{j\ne i}\frac{H_{ij}}{a_j-a_i}=0.
 \tag{I.192}
\]
Since $h_i-h_j=(a_j-a_i)h_ih_j$, (I.190) converts this to
\[
 \sum_{j\ne i}\frac{H_{ij}}{h_i-h_j}=0.
 \tag{I.193}
\]
Indeed the left side of (I.192) is
$h_i^2\sum_{j\ne i}H_{ij}/(h_i-h_j)-h_i\sum_{j\ne i}H_{ij}$.
This derivation remains valid in characteristic two; it
never replaces $Q_H$ by the identically zero alternating
quadratic expression $X^{\mathsf T}HX$ there.

\paragraph{2. A reset-Hessian lemma with these residues.}
Let $h_1,\ldots,h_n$ be distinct elements of a field $K$.
Let $H$ be symmetric with zero diagonal and have the reset
form of I-42:
\[
 H_{ij}=\rho_i\sigma_j\quad(i<j)
 \quad\hbox{if no reset lies strictly between $i,j$},
 \qquad H_{ij}=0\quad\hbox{otherwise}.
 \tag{I.194}
\]
Assume (I.190) and (I.193). Write $t(i)$ for the next
reset strictly after $i$, or $n$ if there is none.
Then
\[
 \begin{array}{ll}
 t(i)-i\le8\ \hbox{for all }i<n&\Longrightarrow H=0
                       \quad\hbox{in every characteristic},\\
 t(i)-i\le10\ \hbox{for all }i<n&\Longrightarrow H=0
                       \quad\hbox{if }\operatorname{char}K\ne2.
 \end{array}
 \tag{I.195}
\]

To prove this, suppose $H\ne0$ and take its first nonzero
row $i$. All earlier columns are also zero. Put $t=t(i)$,
and let $k$ be the first index in $(i,t]$ with
$\sigma_k\ne0$. Then $\rho_i\ne0$. For
$v_j=(1,h_j,h_j^2)^{\mathsf T}$, the row at $i$ and the
row at $k$ give respectively
\[
 \sum_{j=k}^t\sigma_jv_j=0,
 \qquad
 \sum_{j=i}^{k-1}\rho_jv_j=\rho_kv_k.
 \tag{I.196}
\]
Let $J=\{j:i\le j<k,\ \rho_j\ne0\}$ and
$B=\{j:k\le j\le t,\ \sigma_j\ne0\}$, and write
$b=|B|$. The second equation is a nonzero relation because
$\rho_i\ne0$. Any three Vandermonde columns are independent,
so $|J|\ge3$, and $|J|\ge4$ if $\rho_k=0$.

For $j\in J$ one has $\sigma_j=0$, except possibly at
$j=i$ where earlier columns are already zero. Its row has
only the future $B$ contribution. Equation (I.193)
therefore says
\[
 S(h_j)=0,\qquad
 S(T)=\sum_{\ell\in B}\frac{\sigma_\ell}{T-h_\ell}.
 \tag{I.197}
\]
The first relation in (I.196) cancels the first three
terms of the Laurent expansion of $S$ at infinity.
With its denominator $\prod_{\ell\in B}(T-h_\ell)$,
its nonzero numerator thus has degree at most $b-4$.
It is nonzero because each of its distinct simple poles
has nonzero coefficient. Its $|J|$ distinct roots in
(I.197) lie outside its pole set, so
\[
 b\ge |J|+4,\qquad
 t-i\ge(k-i)+(b-1)\ge2|J|+3.
 \tag{I.198}
\]
Since $|J|\ge3$, a nonzero $H$ requires $t-i\ge9$.
This proves the first implication in (I.195).

Outside characteristic two we show $\rho_k=0$. Let
$\ell>k$ be the next index with $\sigma_\ell\ne0$;
it exists because the first relation in (I.196) needs at
least four nonzero terms. For each $k<j<\ell$, its
zero $\sigma_j$ and that same total moment relation make
its three-kernel row equation
\[
 \rho_j\sum_{r=j+1}^{t}\sigma_rv_r
       =-\rho_j\sigma_kv_k=0.
\]
Hence every such $\rho_j$ is zero, with no adjacency
assumption on $k,\ell$. The row at $\ell$ is consequently
\[
 (2\rho_k\sigma_\ell-\rho_\ell\sigma_k)v_k
          -\rho_\ell\sigma_\ell v_\ell=0.
 \tag{I.199}
\]
The two columns are independent. Thus $\rho_\ell=0$ and
then $\rho_k=0$ when $2\ne0$. Now $|J|\ge4$, so (I.198)
requires $t-i\ge11$, proving the second implication.
No sharpness assertion for these new bounds is needed.

\paragraph{3. Application to a rank-preserved scalar word.}
Suppose an actual affine multilinear width-three word has
the preserved interval ranks and nonzero active endpoints
used in I-42, part~3. Its normalized constants are $I_3$
or $\operatorname{diag}(I_2,0)$, with the same transported
endpoint shears. Under a hypothetical two-pole identity,
the evaluated Hessian has exactly (I.194). Part~1 now
supplies both its three kernels and its row residues.
An invertible run of length $L$ makes $t(i)-i$ at most
$L+1$. Part~2 therefore kills $H$ for all runs of length
at most seven, or at most nine outside characteristic two.
The full passive-path argument of I-41 then kills every
coefficient of the scalar word, a contradiction. This is
a conditional rank-preserved-word assertion; the next two
parts provide the required ranks for the stated classes.

\paragraph{4. The constrained connector lemma extends through seven cores.}
Retain the complete anchor classification of I-44. Its
affine-compound case and either pole-good anchor already
work for an arbitrary run length. With the two remaining
anchor forms, the exact determinant is
\[
 F=(\alpha v+\beta A)(dC+e u)-\beta dB_0,
 \quad A=\sum_i p_ix_i,\quad C=\sum_i q_ix_i,
 \quad B_0=\sum_{i<j}p_iq_jx_ix_j,
 \tag{I.200}
\]
where $\alpha\beta de\ne0$. Suppose it annihilates the
one-pole center. Its unique square poles force $p_iq_i=0$.
Put $P=\{i:p_i\ne0\}$ and $Q=\{i:q_i\ne0\}$; these
sets are disjoint. I-44's two quadratic row equations show
that every nonzero row has at least three neighbors.

The row sums at $u$ and $v$, and then at any $q_j\ne0$
and $p_i\ne0$, give respectively
\[
 \begin{gathered}
 \alpha+\beta\sum_i p_i=0,\qquad
 e+d\sum_jq_j=0,\\
 \sum_{i<j}p_i=0\quad(j\in Q),\qquad
 \sum_{j>i}q_j=0\quad(i\in P).
 \end{gathered}
 \tag{I.201}
\]
If $B_0=0$, (I.200) is a product of two nonzero linear
forms and cannot vanish under the one-pole map. Otherwise
choose $i\in P$, $j\in Q$ with $i<j$.
The prefix sum at $j$ in (I.201) contains the nonzero
$p_i$, so there are at least two $P$ positions before $j$.
The row at $j$ has its edge to $v$ and its edges to $P$
positions after $j$; the three-neighbor requirement gives
at least two more $P$ positions after $j$. Hence $|P|\ge4$.
Similarly, the suffix sum at $i$ contains the nonzero
$q_j$ and needs at least two $Q$ positions after $i$,
whereas the three-neighbor condition at $i$ needs two
before it. Thus $|Q|\ge4$.

An annihilating connector with $B_0\ne0$ therefore requires
at least eight internal positions. Every connector with
at most seven full-rank cores survives, in every
characteristic. The same univariate pivots and simultaneous
adjacent-contraction argument of I-44 preserve every
strict internal interval rank for an irreducible homogeneous
translation-invariant minimal word with runs at most seven.
Its endpoint projections also survive by the same pair
dominance. Part~3 now proves scalar nonvanishing for this
entire restricted class in every characteristic.

\paragraph{5. Weight budgets through nine.}
For a homogeneous translation-invariant input, factor and
recompute minimal cores as in I-41. A hypothetical
annihilating factor is irreducible, has nonnegative weight
bounded by the original weight, and all strict internal
ranks are two or three. I-40 bounds its full-rank runs by
its weight $m$.

If $m\le7$, part~4 supplies all ranks and part~3 supplies
the scalar contradiction. For $m\le9$ the only additional
possible run lengths are eight or nine. A length-nine run
is saturated, so its determinant-weight-zero anchor is
covered by I-43's one-invariant-anchor calculation. A
length-eight run has starting determinant weight
$\delta\in\{0,1,2,3\}$ and complementary end weight
$3(m-8)-\delta$. For $m=8$ it is saturated. For $m=9$,
if $\delta=0$ or $3$, one end has weight zero and the same
invariant-anchor argument applies. If $\delta=1$ or $2$,
the three starting weights cannot form a chain
$(q+2,q,q-2)$, since such a chain has determinant weight
$3q$. The affine-compound case of I-44 applies instead.
All ranks and endpoints are therefore preserved even for
$m\le9$ in characteristic two.

When the characteristic is not two, part~3 handles the
resulting maximum run nine. Scalar product closure over
the parameter field completes the factor reduction and
proves (I.187). In characteristic two, the weight-eight
and weight-nine rank statements just proved do not by
themselves prove scalar nonvanishing; the assertion in
(I.187) stops at weight seven there.

\paragraph{6. The finite-variable conclusions.}
Let $M=7$ in arbitrary characteristic, or $M=9$ outside
characteristic two. Consider a width-three multilinear
annihilator on $n\le M+8$ variables. Its lowest nonzero
homogeneous component retains the same coefficient-cut
rank bounds. It is enough to exclude every possible
homogeneous degree.

The known minimum balanced degree and P-19 exclude degrees
at most three. Induct on $d$ from four through
$\lfloor n/2\rfloor$. Every positive translation
coefficient of a degree-$d$ annihilator has smaller degree.
If any were nonzero, the highest such coefficient would
inherit width at most three by the leading-parameter
argument of I-40, contradicting the induction hypothesis.
Thus the degree-$d$ annihilator is translation-invariant.
Its weight satisfies
$0\le n-2d\le n-8\le M$, so (I.187) excludes it.
Multilinear reversal excludes all complementary degrees
above $n/2$. This proves the fifteen- and seventeen-variable
assertions, including inhomogeneous inputs.

The denominator-safe interpolation of I-40 is unchanged:
the two selector degrees are at most $n$ and the total
degree in the three amplitude variables is at most $D$.
It gives (I.188) with precisely the stated field scope.
\hfill$\square$

\paragraph{Scope of the improvement.}
The additional row equations are a consequence of a fixed
polynomial two-pole identity, not a condition inferred from
an arbitrary symmetric matrix. They explain why I-42's
abstract seven-node characteristic-two example does not
obstruct the new scalar theorem. No sharpness or polynomial
integrability is asserted for the new residue bounds.

At larger weights, the two-remaining-anchor determinant
(I.200) has at least eight internal positions in any possible
one-pole obstruction. Higher-order Taylor constraints may
also be relevant beyond the present Hessian thresholds.
Those questions, the full balanced width-three theorem,
and polynomial dependence on growing ROABP width remain
open here. No historical priority claim or numerical
verification is asserted.
